\newpage

%~
%\vfill
\HalfPageImage{forprint/140_folio_42_recto_B_OF_KP_III.jpg}{\krnt, folio 42\,\RE.}
%\vskip -10pt

%\Rline{\krnt, folio 42\,\RE.}
%\bigskip
%\vfill

\HalfPageImage{forprint/141_folio_42_verso_E_OF_KP_III.jpg}{\krnt, folio 42\,\VE.}
%\vskip -10pt

%\Rline{\krnt, folio 42\,\VE.}
%\vfill

\newpage
%~\vfill

\begin{HalfText}
\begin{sans}
\begin{center}
\booktitle{[goraṇṭakapaddhatiḥ]}
\medskip

{\LARGE\bfseries śrīgaṇeśāya namaḥ}
\end{center}

prathamābhyāsaprārambhe vighna ālasyamucyate\!। 
ālasyaṃ bhakṣayedyogī dantānkaṭakaṭāyya\endnote{dantān~] \textit{emend}., dantā+ M. kaṭakaṭāyya~] M(pc), kaṭakaṭā+ya M(ac).} ca\!॥1\!॥\\
mandībhūte 'tha\endnote{'tha~] \textit{avagraha added}, tha M.} cālasye vedanāṅge prajāyate\!। 
vedanābhyāsayogena āsaneṣu vilīyate\!॥2\!॥\\
āsanaṃ tu dṛḍhīkāryaṃ vedanāśāntihetave\!। 
abhyāsakṛtavaiṣamyam abhyāsena tu nirdahet\endnote{nirdahet~] M(pc), niṃdahet M (ac).}\!॥3\!॥\\
na ca vaidyaṃ samāśritya roganirdahane kṣamaḥ\!। 
āsane tu dṛḍhībhūte vāyumabhyāsamabhyaset\!॥4\!॥\\
svedaparyantamabhyāsaḥ kartavyaḥ saṃbhrameṇa tu\!। 
śarīre gharmaklinne ca bhasmamardanamiṣyate\!॥5\!॥\\
bhasmaśayyāṃ tataḥ kuryāj jantupīḍanamārgikīm\!। 
tatraiva nidrāṃ kurvīta dehasauṣṭhavakāmyayā\!। 
paridhāryaṃ sūkṣmavastraṃ rāgasyūtyādivarjitam\endnote{°syūty°~] \textit{diagnostic conj.} \textsc{Goodall}, °nīty° M. This conjecture, which retains one ligature of the transmitted reading and fits the metre, is based on the idea that a fine cloth (\textit{sūkṣmavastra}) was probably considered to be smooth (e.g., \textit{Śabdakalpadruma} (s.v., 39911): \textit{sūkṣmavastra} = \textit{ślakṣṇavasana}), in which case it would not have stitching, hems, etc.}\!॥6\!॥\\
śarīre kampa utpanne cābhyāso madhyatāṃ vrajet\!। 
kampe jīrṇe nāḍiśuddhiḥ\endnote{°śuddhiḥ~] \textit{corr.}, °śudhiḥ M.} syāttadaiva\endnote{tadaiva~] \textit{emend}. \textsc{Goodall}, tadeva M.} hi yoginaḥ\!॥7\!॥\\
āpādamastakaṃ vāyuḥ krīḍate lāghavena tu\!। 
netre\endnote{netre~] M(pc), nehetre M(ac).} nirmalatāṃ yātaḥ śarīre lāghavaṃ tathā\!॥8\!॥\\
danteṣu dārḍhyaṃ bhavati nakheṣvāraktatā tataḥ\!। 
pāmādirogā naśyanti kuṣṭhādīnāṃ ca nāśanam\!॥9\!॥\\
arśogulmādināśena\endnote{arśo°~] \textit{emend}., ārṣo° M.} śarīre pāṭavaṃ bahu\!। 
śaitye vikāre mārīcaḥ paitte\endnote{paitte~] \textit{corr.}, paite M.} brāhmī\endnote{brāhmī~] \textit{emend}., brāṃhmī M.} praśasyate\!॥10\!॥\\
vāte jyotiṣmatī proktā mahatā kākinā khalu\!। 
taṇḍulā godhumā uktāḥ\endnote{uktāḥ~] \textit{corr.}, uktāṣ M(pc), uktā+ M(pc).} ṣaṣṭikādyāśca yogibhiḥ\!॥11\!॥\\
teṣāṃ tu bhakṣaṇaṃ śastaṃ dṛḍhe 'bhyāse\endnote{'bhyāse~] \textit{avagraha added}, bhyāse M.} tu na kṣatiḥ\!। 
ete guṇāśca draṣṭavyā vajrolīpāramicchunā\endnote{icchunā~] \textit{corr.}, ichunā M.}\!॥12\!॥\\
mudrācihnāni\endnote{°cihnāni~] \textit{corr}, °cinhāni M.} dhāryāṇi helayā vātha līlayā\!। 
kṛṣṇājinadvayaṃ sāṅgaṃ khuraśṛṅgādiśobhitam\endnote{°śṛṅgādi°~] \textit{conj.} \textsc{Goodall}, °ādiśṛṃ° M.}\!॥13\!॥\\
vyāghrāmbaraṃ tathā caikaṃ kanthākambalasambhavam\!। 
kamaṇḍuluṃ\endnote{The standard spelling of this word is \textit{kamaṇḍalu}. However, while not attested in the dictionary, the spelling \textit{kamaṇḍulu} occurs elsewhere, including the \textit{Śrītattvanidhi} (\textit{kautukanidhi} 7.8).} khaḍgajaṃ ca tantrīvādyaṃ ca pustakam\endnote{pustakam~] \textit{emend.} \textsc{Mallinson}, mastakam M}\!॥14\!॥\\
mastake ca jaṭā\insmark{0pt}{0pt}{jū}ṭaṃ kāntasaṃskāravarjitam\endnote{kānta°~] \textit{conj}., kanta° M.}\!।
\end{sans}
\end{HalfText}
%\vfill

\begin{HalfText}
\begin{sans}
śvetaṃ\endnote{śvetaṃ~] \textit{corr.}, svetaṃ M.} ca bhasma saṃdhāryam\endnote{°dhāryam~] \textit{corr.}, °dhāryaṃm M.} aṅge gomayanirmitam\!॥15\!॥\\
kaupīnaṃ sūkṣmavastraṃ ca daṇḍaṃ ca veṇusaṃbhavam\!। 
dārumaye\endnote{°maye~] \textit{emend}., °mayau M.} pāduke ca dhārye\endnote{dhārye~] \textit{emend}., dhāryā M.} vai siddhimicchatā\endnote{siddhim~] M(pc), siddhide M(ac). icchatā~] \textit{corr.}, ichatā M.}\!॥16\!॥\\
vāmabhāgaviśuddhyarthaṃ śiṣyerṣyājananāya ca\!। 
suvarṇanigaḍaṃ dhṛtvā cāṅguṣṭhe valayadvayam\!। 
mudrāsiddhastu\endnote{°siddhas~] \textit{emend}, ° siddhis M.} yaḥ śiṣyaḥ gṛhṇīyānnigaḍaṃ\endnote{gṛhṇīyān~] \textit{corr.},  gṛṇhīyān M.} guroḥ\!॥17\!॥\\
suvarṇasyeṣṭakāṃ dhṛtvā koṇairdvādaśabhiryutām\endnote{dvādaśabhir~] \textit{emend.}, dvādabhir M. yutām~] \textit{emend}., yutā M.}\!। 
tāvadbhiḥ koṣṭhakairyuktā akṣaraiśca suśobhitā\!॥18\!॥\\
gandhairaṣṭabhirāliptābhyaktā yakṣasya kardamaiḥ\!। 
ācchādayedvedikāyāṃ\endnote{ācchādayed~] \textit{corr.}, āchādayed M.} gandhapuṣpaiḥ prapūjayet\!॥19\!॥\\
tatsaṃnidhau japaṃ kuryādvādaśākṣaravidyayā\!। 
śatrunāśo mitravṛddhiḥ dhanavṛddhiśca jāyate\!॥20\!॥\\
rājānaśca vaśe tasya bhaveyur nātra saṃśayaḥ\!। 
vāsudevaśca santuṣṭo dadetsāyujyatāmapi\!॥21\!॥\\
japāddviguṇato\endnote{japād~] \textit{emend}., japā M.} homo homāddviguṇatarpaṇam\endnote{homād~] \textit{emend}., homā M.}\!। 
mahāpūjā ca kartavyā mantrairvaidikatāntrikaiḥ\!॥22\!॥\\
bhojanaṃ brāhmaṇebhyaśca dātavyaṃ siddhimicchatā\endnote{icchatā~] \textit{corr.}, ichatā M.}\!। 
viduṣaḥ\endnote{viduṣaḥ~] \textit{emend.} \textsc{Goodall}, viduṣā M.} pūjayedbhaktyā\endnote{pūjayed~] \textit{conj.}, pūja pūjayed M (\textit{dittography}).} sarvasiddheśca\endnote{°siddheś~] \textit{emend}., °siddhiś M.} kāraṇāt\!॥23\!॥\\
iti mudrā yoge guptābhyāsataḥ prakaṭīkṛtā\!। 
goraṇṭakenātmakṛte likhitā śiṣyabodhane\!॥24\!॥\\
iti goraṇṭakapaddhatiruddhṛtā\endnote{°paddhatir~] \textit{emend}. \textsc{Sanderson}, °paddhatim M.} yairadhigamya dharmyamarmaśikṣāḥ\!। 
te 'mutra ca kāryadakṣā\endnote{ca kārya°~] \textit{emend}. \textsc{Goodall}, cākārya° M.} rakṣā teṣāṃ prabhubhiḥ kṛtā samakṣā\!॥25\!॥\\
iti kuraṇṭakapaddhatiḥ samāptā॥
\end{sans}
\end{HalfText}
%\vfill

\newpage

~\vskip 1cm
\begin{eng}
\begin{center}
\booktitle{[Goraṇṭaka’s Manual on Yoga]}
\medskip

{\LARGE\bfseries Homage to Gaṇeśa}
\end{center}
\end{eng}


\begin{multicols}{2}
\begin{eng}
\Section{[The Practice]}

\begin{enumerate}%[nosep]
\item[1.] When first undertaking the practice, a [yogi‘s main] obstacle is said to be laziness. Gritting his teeth, the yogi should devour laziness.\endnote{The expression \textit{dantān kaṭakaṭāyya} (literally, “gnashing the teeth”) occurs elsewhere in Sanskrit literature, where it is sometimes used to describe a person who is angry or in a rage (e.g., \textit{Rāmāyaṇa} 7.61.2, \textit{Mahābhārata} 5.159.6 [12.25], etc.). In the current context, it appears to be a metonymic expression for determined effort, hence our translation “gritting the teeth.”}
\item[2.] When laziness subsides, pain then arises in the body. The pain in the \textit{āsana}s dissolves through practice.
\item[3.] [Each] \textit{āsana} should be made firm to alleviate pain. An imbalance caused by practice is burnt up by the practice.
\item[4.] And without resorting to a doctor, [a yogi] should be capable of burning up diseases.  When \textit{āsana} has become steady, one should do the breath practice.
\item[5.] The [breath] practice should be done vigorously until sweat arises.\endnote{Various other texts mention that sweat (\textit{sveda}) and shaking (\textit{kampa}, mentioned in verse 7) are signs of the lowest and intermediate stages of \textit{prāṇāyāma}, respectively. Examples include \textit{Haṭhapradīpikā} 2.12 and \textit{Vivekamārtaṇḍa} 87.} And when the body is moistened by the heat, rubbing it with ashes is prescribed.\endnote{Cf.\ \textit{Haṭhapradīpikā} 2.13. The \textit{Haṭhapradīpikā} states that one should rub sweat into the body, but it does not prescribe the use of ash for this.}
\item[6.]  One should then make a bed of ashes, which is a way to suppress insects, [and] go to sleep on it with the desire for physical wellbeing. One should wrap around [oneself] a fine cloth, free of dye, stitching, and so on.
\item[7.] When trembling has arisen in the body, the [breath] practice progresses to an intermediary stage.  When the trembling has diminished, only at that time does the yogi‘s channels become pure.
\end{enumerate}
\bigskip
~\phantom{ii}\\
~\phantom{ii}\\
~\phantom{ii}\\
\Section{[Health Benefits and Remedies]}

\begin{enumerate}%[nosep]
\item[8.] [Then] the breath playfully moves with lightness as far as the feet and the head. The eyes become pure and there is lightness in the body.
\item[9.] There is strength in the teeth and then a slight reddishness in the nails. Ailments, such as scabby eruptions and so on, and [other] skin diseases disappear.\endnote{The skin disease \textit{pāmā} is usually classified as a minor variety of \textit{kuṣṭha}, which is a broad category of skin diseases (e.g., \textit{Carakasaṃhitā, Cikitsāsthāna} 7.21–26). It is often translated as scabies (Meulenbeld, vol 1b, 1999: 109, note 199) but the general sense of scabies as a scabby eruption might be intended rather than the specific disease caused by the itch mite, \textit{Sarcoptes scabiei}.}
\item[10.] With the disappearance of piles, swelling and so on, optimal health arises in the body.\endnote{The compound \textit{aṅgapāṭava} is used in an early version of the \textit{Haṭhapradīpikā} (1.17) to mean optimal functioning of the body. In the \textit{Śivasaṃhitā} (2.35),  \textit{śarīrapāṭava} is mentioned as one of the benefits bestowed by digestive fire (\textit{vaiśvānarāgni}).} Pepper is recommended for cold ailments and Brāhmī for those of bile.
\item[11.] As is well known, [taking] the Black Oil plant for wind ailments has been taught by the great Kāki.\endnote{\textit{Jyotiṣmatī} is usually identified as \textit{Celastrus paniculatus} \textsc{Willd} (e.g., Singh and Chunekar 1999: 171), which is commonly known as the black oil plant, climbing staff tree, and intellect tree. On the possible identity of Kāki, see the introduction.} Yogis recommend rice, wheat, sixty-day rice, and the like.\endnote{Both \textit{taṇḍula} and \textit{godhuma} are mentioned as wholesome foods in the \textit{Kapālakurāṇṭaka}, and \textit{ṣaṣṭika} is recommended in \textit{Haṭhapradīpikā} 1.62.}
\item[12.] Consuming these [herbs and foods] is recommended. When the practice is steady, there is no harm [from taking them]. They are to be seen as beneficial for one who desires to reach the farther shore [of transmigration] by [the practice of] \textit{vajrolī}.
\end{enumerate}
%\newpage
\bigskip

\Section{[The Yogi]}

\begin{enumerate}%[nosep]
\item[13–14.] Marks and emblems should be worn mockingly or playfully.\endnote{We do not know of a parallel in another text for the notion of yogis wearing marks and emblems mockingly or playfully. The closest idea may be that expressed in \textit{Amanaska} 2.34–36, which states that adopting sectarian emblems (\textit{liṅga}) is not useful in any situation for one who is self-realised and detached.} [The yogi should carry] two complete black antelope skins adorned with horns, hooves, and so on; a tiger-skin garment; one made from rags and woollen blankets; a water pot made from the horn of a rhinoceros; a stringed instrument; and a book.\endnote{Since it is unlikely that a yogi would be carrying a ‘head,’ the manuscript’s reading of \textit{mastakam} at the end of 14d is likely a dittographic error as the next verse begins with \textit{mastake}.}
\item[15.] Also, [the yogi should wear] a twisted mass of matted locks on his head, free from beautiful adornments. And white ash made of cow dung should be worn on the body.
\item[16.] [He should have] a loin cloth, a [piece of] fine cloth,\endnote{Cf.~verse 6.} and a staff made of bamboo. [The yogi] desirous of success should wear sandals made of wood.
\end{enumerate}
\bigskip

\Section{[Mantra Recitation with a Golden Brick]}

\begin{enumerate}%[nosep]
\item[17.] For purifying the left side\endnote{We are not sure what is meant by the ‘left side’ (\textit{vāmabhāga}) here. It is unclear whether this refers to someone positioned on the student‘s left side, such as a consort; the left side of the body; or something within the left side of the body, such as the \textit{iḍā} channel.} and to arouse envy in [other] students, a student who has been perfected by the \textit{mudrā} should wear a gold chain and two rings on the thumb, and should receive a chain from the guru.
\item[18.] He should hold a golden sacrificial brick with twelve corners. It should be endowed with as many facets and nicely decorated with a syllable [on each facet].\endnote{It seems probable that each facet contained one syllable of the twelve-syllable Vāsudeva mantra, which is attested in full in the \textit{Kapālakuraṇṭakapaddhati} and alluded to below.}
\item[19.] [The brick] is smeared with eight fragrant substances [and] anointed with perfumed pastes [called \textit{yakṣakardama}]. [The student] should cover it and worship it on an altar with fragrant flowers. 
\item[20.] Near it, he should recite the twelve-syllable mantra. It destroys enemies and increases friends and wealth.
\item[21.] And kings would undoubtedly be under his control. Satisfied, Vāsudeva also bestows [liberation in the form of being in] his presence. 
\item[22.] Fire sacrifice is twice as effective as mantra recitation, and libation is twice as effective as fire sacrifice. The [most] exalted worship should be performed with vedic and tantric mantras.
\item[23.] A person desirous of success should give food to Brahmins. With devotion one should honour the wise, for this is the cause of all success.
\end{enumerate}
\bigskip

\Section{[Conclusion]}

\begin{enumerate}%[nosep]
\item[24.] Thus, the \textit{mudrā}, which is secret in yoga and has been revealed through practice, has been written down by Goraṇṭaka for his own sake and for awakening students.
\item[25.] Those who have understood the teachings on the secret points of virtue and have extracted the manual of Goraṇṭaka are skilled in accomplishing whatever needs to be done and have the visible protection afforded by their lords, even in the next world.
\end{enumerate}
Thus, the manual of Kuraṇṭaka is complete.
\end{eng}
\end{multicols}

\newpage

\begin{kan}
~\vskip 1cm
%ಕೋರಂಟಕ ಪದ್ಧತಿ - ತೃತೀಯ ಅಧ್ಯಾಯ
\begin{center}
\booktitle{ಗೊರಂಟಕನ ಯೋಗ ಕೈಪಿಡಿ}
\addcontentsline{tocka}{chapter}{\kantoc{ಗೊರಂಟಕನ ಯೋಗ ಕೈಪಿಡಿ}}
\medskip

{\LARGE\bfseries ಶ್ರೀಗಣೇಶಾಯ ನಮಃ}

{\LARGE\bfseries ಗಣೇಶನಿಗೆ ನಮನಗಳು.}
\end{center}
\end{kan}

\begin{multicols}{2}
\begin{kan}
\Section{ಅಭ್ಯಾಸ ಕ್ರಮ}

\begin{enumerate}%[nosep]
\item[1.] ಅಭ್ಯಾಸವನ್ನು ಪ್ರಾರಂಭಿಸಿದಾಗ, ಯೋಗಿಗೆ ಎದುರಾಗುವ ಮೊದಲ ಮತ್ತು ಮುಖ್ಯ ಅಡೆತಡೆ ಎಂದರೆ ಅದು ‘ಆಲಸ್ಯ’. ಯೋಗಿಯು ತನ್ನ ಹಲ್ಲುಗಳನ್ನು ಕಚ್ಚಿ (ಛಲದಿಂದ), ಆ ಆಲಸ್ಯವನ್ನು ನುಂಗಿ ಹಾಕಬೇಕು ಅಂದರೆ ಅದನ್ನು ಸಂಪೂರ್ಣವಾಗಿ ಜಯಿಸಬೇಕು.
\item[2.] ಆಲಸ್ಯವು ಕಡಿಮೆಯಾದಾಗ, ದೇಹದಲ್ಲಿ ನೋವು ಕಾಣಿಸಿಕೊಳ್ಳುತ್ತದೆ. ಆಸನಗಳನ್ನು ಮಾಡುವಾಗ ಉಂಟಾಗುವ ಈ ನೋವು ನಿರಂತರ ಅಭ್ಯಾಸದ ಮೂಲಕ ಕ್ರಮೇಣ ಕರಗಿಹೋಗುತ್ತದೆ.
\item[3.] ನೋವನ್ನು ನಿವಾರಿಸಲು ಪ್ರತಿಯೊಂದು ಆಸನವನ್ನೂ ದೃಢವಾಗಿ ಮಾಡಬೇಕು. ಅಭ್ಯಾಸದಿಂದ ಉಂಟಾದ ದೈಹಿಕ ಅಸಮತೋಲನಗಳು ಅದೇ ಅಭ್ಯಾಸದ ತೀವ್ರತೆಯಿಂದ ಸುಟ್ಟು ಹೋಗುತ್ತವೆ.
\item[4.] ವೈದ್ಯರ ಸಹಾಯವಿಲ್ಲದೆಯೇ ಕಾಯಿಲೆಗಳನ್ನು ಗುಣಪಡಿಸುವ ಸಾಮರ್ಥ್ಯವನ್ನು ಯೋಗಿಯು ಬೆಳೆಸಿಕೊಳ್ಳಬೇಕು. ಆಸನವು ಸ್ಥಿರವಾದ ನಂತರ, ಪ್ರಾಣಾಯಾಮ ಅಥವಾ ಉಸಿರಾಟದ ಅಭ್ಯಾಸವನ್ನು ಮಾಡಬೇಕು.
\item[5.] ಬೆವರು ಬರುವವರೆಗೆ ಉಸಿರಾಟದ ಅಭ್ಯಾಸವನ್ನು ಅತ್ಯಂತ ವೇಗವಾಗಿ ಅಥವಾ ಬಲವಾಗಿ ಮಾಡಬೇಕು. ದೇಹವು ಆ ಉಷ್ಣತೆಯಿಂದ (ಬೆವರಿನಿಂದ) ತೇವವಾದಾಗ, ದೇಹಕ್ಕೆ ಭಸ್ಮವನ್ನು ಹಚ್ಚಿಕೊಳ್ಳುವುದು ವಿಹಿತವಾಗಿದೆ.
\item[6.] ನಂತರ ಕೀಟಗಳ ಬಾಧೆಯನ್ನು ತಪ್ಪಿಸಲು ಭಸ್ಮವನ್ನು ಹಾಸಿಗೆಯಂತೆ ಹರಡಿ, ಶಾರೀರಿಕ ಆರೋಗ್ಯದ ಹಿತದೃಷ್ಟಿಯಿಂದ ಅದರ ಮೇಲೆ ಮಲಗಬೇಕು. ಬಣ್ಣ ಅಥವಾ ಹೊಲಿಗೆಯಿಲ್ಲದ ಒಂದು ಶುಭ್ರವಾದ ಮೃದು ಬಟ್ಟೆಯನ್ನು ಮೈಮೇಲೆ ಹೊದ್ದುಕೊಳ್ಳಬೇಕು.
\item[7.] ಉಸಿರಾಟದ ಅಭ್ಯಾಸ ಮಾಡುತ್ತಿರುವಾಗ ಯಾವಾಗ ದೇಹದಲ್ಲಿ ನಡುಕ ಉಂಟಾಗುತ್ತದೆಯೋ, ಆಗ ಯೋಗಿಯ ಸಾಧನೆಯು ಮಧ್ಯಮ ಹಂತಕ್ಕೆ ತಲುಪಿದ ಸಂಕೇತ ಎಂದು ತಿಳಿಯಬೇಕು. ಆ ನಡುಕವು ನಿಂತಾಗ ಮಾತ್ರ ಯೋಗಿಯ ನಾಡಿಗಳು ಸಂಪೂರ್ಣವಾಗಿ ಶುದ್ಧವಾಗಿವೆ ಎಂದರ್ಥ.
\end{enumerate}
\bigskip

\Section{ಆರೋಗ್ಯದ ಪ್ರಯೋಜನಗಳು ಮತ್ತು ಪರಿಹಾರಗಳು}

\begin{enumerate}%[nosep]
\item[8.] ನಾಡಿಗಳು ಶುದ್ಧವಾದ ನಂತರ ಉಸಿರಾಟವು ಅತ್ಯಂತ ಹಗುರವಾಗಿ ಪಾದಗಳಿಂದ ತಲೆಯವರೆಗೆ ಲೀಲಾಜಾಲವಾಗಿ ಸಂಚರಿಸುತ್ತದೆ. ಕಣ್ಣುಗಳು ನಿರ್ಮಲವಾಗುತ್ತವೆ ಮತ್ತು ಶರೀರದಲ್ಲಿ ಲಘುತ್ವ (ಹಗುರವಾದ ಭಾವನೆ) ಉಂಟಾಗುತ್ತದೆ.
\item[9.] ಹಲ್ಲುಗಳಲ್ಲಿ ಬಲವುಂಟಾಗುತ್ತದೆ ಮತ್ತು ಉಗುರುಗಳಲ್ಲಿ ಸ್ವಲ್ಪ ಕೆಂಪು ಬಣ್ಣ ಕಾಣಿಸಿಕೊಳ್ಳುತ್ತದೆ. ಕಜ್ಜಿ ಅಥವಾ ತುರಿಕೆಯಂತಹ ಚರ್ಮದ ಕಾಯಿಲೆಗಳು ಮತ್ತು ಇತರ ವ್ಯಾಧಿಗಳು ಮಾಯವಾಗುತ್ತವೆ.
\item[10.] ಮೂಲವ್ಯಾಧಿ, ಬಾವು ಮುಂತಾದ ತೊಂದರೆಗಳು ಗುಣವಾಗುವುದರೊಂದಿಗೆ ಶರೀರದಲ್ಲಿ ಉತ್ತಮ ಆರೋಗ್ಯವು ನೆಲೆಸುತ್ತದೆ.
\item[11.] ವಾಯು ದೋಷದ ನಿವಾರಣೆಗಾಗಿ ಜ್ಯೋತಿಷ್ಮತಿ ಗಿಡವನ್ನು ಬಳಸುವ ಕ್ರಮವನ್ನು ಮಹಾನ್ ಕಾಕಿಯು ಬೋಧಿಸಿದ್ದಾನೆ. ಯೋಗಿಗಳು ಅಕ್ಕಿ, ಗೋಧಿ ಮತ್ತು ಅರವತ್ತು ದಿನದ ಅಕ್ಕಿ ಮುಂತಾದ ಆಹಾರಗಳನ್ನು ಶಿಫಾರಸು ಮಾಡುತ್ತಾರೆ.
\item[12.] ಇಂತಹ ಗಿಡಮೂಲಿಕೆಗಳು ಮತ್ತು ಆಹಾರಗಳನ್ನು ಸೇವಿಸುವುದು ಒಳ್ಳೆಯದು. ಯೋಗಾಭ್ಯಾಸವು ಸ್ಥಿರವಾಗಿದ್ದಾಗ, ಇವುಗಳ ಸೇವನೆಯಿಂದ ಯಾವುದೇ ಹಾನಿಯಾಗುವುದಿಲ್ಲ. ವಜ್ರೋಲಿಯ ಮೂಲಕ ಸಂಸಾರ ಸಾಗರದ ದಡ ಸೇರಲು ಬಯಸುವವರಿಗೆ ಇವು ಪ್ರಯೋಜನಕಾರಿ ಎಂದು ಪರಿಗಣಿಸಬೇಕು.
\end{enumerate}
\bigskip

\Section{ಯೋಗಿಯ ಲಕ್ಷಣಗಳು}
 
\begin{enumerate}%[nosep]
\item[13–14.] ಯೋಗಿಯು ಲೋಕವನ್ನು ಅಣಕಿಸುವಂತೆ ಅಥವಾ ಲೀಲಾಜಾಲವಾಗಿ ಚಿಹ್ನೆಗಳನ್ನು ಮತ್ತು ಮುದ್ರೆಗಳನ್ನು ಧರಿಸಬೇಕು. ಯೋಗಿಯು ಕೊಂಬುಗಳು ಮತ್ತು ಗೊರಸುಗಳಿಂದ ಅಲಂಕರಿಸಲ್ಪಟ್ಟ ಎರಡು ಪೂರ್ಣ ಪ್ರಮಾಣದ ಕೃಷ್ಣಮೃಗದ ಚರ್ಮಗಳು; ಹುಲಿಚರ್ಮದ ವಸ್ತ್ರ; ಚಿಂದಿ ಬಟ್ಟೆಗಳು ಮತ್ತು ಉಣ್ಣೆಯ ಕಂಬಳಿಗಳಿಂದ ಮಾಡಿದ ಉಡುಪು; ಖಡ್ಗಮೃಗದ ಕೊಂಬಿನಿಂದ ಮಾಡಿದ ಕಮಂಡಲ (ನೀರಿನ ಪಾತ್ರೆ); ಒಂದು ತಂತೀವಾದ್ಯ ಮತ್ತು ಒಂದು ಪುಸ್ತಕ ಇವುಗಳನ್ನು ಹೊಂದಿರಬೇಕು.
\item[15.] ಅಲ್ಲದೆ, ಯೋಗಿಯು ತನ್ನ ತಲೆಯ ಮೇಲೆ ಯಾವುದೇ ಸುಂದರ ಆಭರಣಗಳಿಲ್ಲದ ಜಟೆಯನ್ನು (ಹೆಣೆದ ಕೂದಲು) ಧರಿಸಿರಬೇಕು. ಮತ್ತು ಗೋಮಯದಿಂದ (ಆಕಳ ಸಗಣಿ) ಮಾಡಿದ ಬಿಳಿ ಭಸ್ಮವನ್ನು ಶರೀರದಾದ್ಯಂತ ಧರಿಸಿರಬೇಕು.
\item[16.] ಅವನು ಒಂದು ಕೌಪೀನ, ಒಂದು ಶುಭ್ರವಾದ ಬಟ್ಟೆ ಮತ್ತು ಬಿದಿರಿನ ದಂಡವನ್ನು ಹೊಂದಿರಬೇಕು. ಸಿದ್ಧಿಯನ್ನು ಬಯಸುವ ಯೋಗಿಯು ಮರದಿಂದ ಮಾಡಿದ ಪಾದರಕ್ಷೆಗಳನ್ನು (ಮರದ ಪಾದುಕೆಗಳು) ಧರಿಸಬೇಕು.
\end{enumerate}
\bigskip

\Section{ಚಿನ್ನದ ಇಟ್ಟಿಗೆಯೊಂದಿಗೆ ಮಂತ್ರ ಪಠಣೆ}

\begin{enumerate}%[nosep]
\item[17.] ಎಡಭಾಗವನ್ನು ಶುದ್ಧೀಕರಿಸಲು ಮತ್ತು ಇತರ ವಿದ್ಯಾರ್ಥಿಗಳಲ್ಲಿ ಆಸಕ್ತಿಯನ್ನು ಹುಟ್ಟಿಸಲು, ಮುದ್ರೆಯಲ್ಲಿ ಸಿದ್ಧಿಯನ್ನು ಪಡೆದ ಶಿಷ್ಯನು ಚಿನ್ನದ ಸರ ಮತ್ತು ಹೆಬ್ಬೆರಳಿಗೆ ಎರಡು ಉಂಗುರಗಳನ್ನು ಧರಿಸಬೇಕು. ಅಲ್ಲದೆ, ಅವನು ಗುರುವಿನಿಂದ ಒಂದು ಸರವನ್ನು ಪಡೆಯಬೇಕು.
\item[18]. ಅವನು ಹನ್ನೆರಡು ಮೂಲೆಗಳಿರುವ ಚಿನ್ನದ ಯಜ್ಞದ ಇಟ್ಟಿಗೆಯನ್ನು ಹಿಡಿದಿರಬೇಕು. ಅದು ಹನ್ನೆರಡು ಮುಖಗಳನ್ನು ಹೊಂದಿರಬೇಕು ಮತ್ತು ಪ್ರತಿಯೊಂದು ಮುಖದ ಮೇಲೆ ಒಂದು ಅಕ್ಷರವನ್ನು (ಮಂತ್ರದ ಅಕ್ಷರ) ಸುಂದರವಾಗಿ ಕೆತ್ತಿರಬೇಕು.
\item[19.] ಆ ಇಟ್ಟಿಗೆಗೆ ಎಂಟು ಸುಗಂಧ ದ್ರವ್ಯಗಳನ್ನು ಹಚ್ಚಬೇಕು ಮತ್ತು ‘ಯಕ್ಷಕರ್ದಮ’ ಎಂಬ ಸುಗಂಧ ಲೇಪನದಿಂದ ಪೂಜಿಸಬೇಕು. ಶಿಷ್ಯನು ಅದನ್ನು ಮುಚ್ಚಿಟ್ಟು, ಹೂವುಗಳಿಂದ ಅಲಂಕರಿಸಿದ ಪೀಠದ ಮೇಲೆ ಪೂಜಿಸಬೇಕು.
\item[20.] ಅದರ ಮುಂದೆ ಕುಳಿತು ದ್ವಾದಶಾಕ್ಷರಿ ಮಂತ್ರವನ್ನು ಜಪಿಸಬೇಕು. ಇದು ಶತ್ರುಗಳನ್ನು ನಾಶಪಡಿಸುತ್ತದೆ ಮತ್ತು ಮಿತ್ರರನ್ನು ಹಾಗೂ ಸಂಪತ್ತನ್ನು ವೃದ್ಧಿಸುತ್ತದೆ.
\item[21.] ಇದರಿಂದ ರಾಜರು ನಿಸ್ಸಂದೇಹವಾಗಿ ಯೋಗಿಯ ವಶವಾಗುತ್ತಾರೆ. ಪ್ರಸನ್ನನಾದ ವಾಸುದೇವನು ತನ್ನ ಸಾನ್ನಿಧ್ಯದ ರೂಪದಲ್ಲಿ ಮುಕ್ತಿಯನ್ನು ಕರುಣಿಸುತ್ತಾನೆ.
\item[22.] ಮಂತ್ರ ಪಠಣೆಗಿಂತ ಹೋಮವು ಎರಡರಷ್ಟು ಹೆಚ್ಚು ಪರಿಣಾಮಕಾರಿ, ಮತ್ತು ಹೋಮಕ್ಕಿಂತ ತರ್ಪಣವು ಎರಡರಷ್ಟು ಶ್ರೇಷ್ಠ. ವೇದ ಮತ್ತು ತಂತ್ರೋಕ್ತ ಮಂತ್ರಗಳ ಮೂಲಕ ಅತ್ಯುನ್ನತವಾದ ಪೂಜೆಯನ್ನು ಸಲ್ಲಿಸಬೇಕು.
\item[23.] ಸಿದ್ಧಿಯನ್ನು ಬಯಸುವವನು ಬ್ರಾಹ್ಮಣರಿಗೆ ಅನ್ನದಾನ ಮಾಡಬೇಕು. ಭಕ್ತಿಯಿಂದ ಜ್ಞಾನಿಗಳನ್ನು ಗೌರವಿಸಬೇಕು, ಏಕೆಂದರೆ ಇದೇ ಎಲ್ಲಾ ಯಶಸ್ಸಿಗೆ ಮೂಲಕಾರಣ.
\end{enumerate}
\bigskip

\Section{ಸಮಾಪ್ತಿ}

\begin{enumerate}%[nosep]
\item[24.] ಈ ರೀತಿಯಾಗಿ, ಯೋಗದಲ್ಲಿ ಅತ್ಯಂತ ರಹಸ್ಯವಾಗಿರುವ ಮತ್ತು ಕಠಿಣ ಅಭ್ಯಾಸದ ಮೂಲಕ ಮಾತ್ರ ಅರಿವಿಗೆ ಬರುವ ಈ ಮುದ್ರೆಯನ್ನು, ಗೊರಂಟಕನು ತನ್ನ ಸ್ವಂತ ತಿಳುವಳಿಕೆಗಾಗಿ ಮತ್ತು ಜಾಗೃತಗೊಳ್ಳುತ್ತಿರುವ (ಕಲಿಯುತ್ತಿರುವ) ವಿದ್ಯಾರ್ಥಿಗಳಿಗಾಗಿ ಬರೆದಿದ್ದಾನೆ.
\item[25.] ಯೋಗದ ಪುಣ್ಯದಾಯಕ ರಹಸ್ಯಗಳನ್ನು ಯಾರು ಅರ್ಥಮಾಡಿಕೊಂಡಿದ್ದಾರೋ ಮತ್ತು ಗೊರಂಟಕನ ಈ ಕೈಪಿಡಿಯ ಸಾರವನ್ನು ಯಾರು ಗ್ರಹಿಸಿದ್ದಾರೋ, ಅವರು ಮಾಡಬೇಕಾದ ಕೆಲಸಗಳನ್ನು ಸಾಧಿಸುವುದರಲ್ಲಿ ನಿಪುಣರಾಗುತ್ತಾರೆ. ಅಷ್ಟೇ ಅಲ್ಲದೆ, ಅವರಿಗೆ ಈ ಲೋಕದಲ್ಲಿ ಮತ್ತು ಪರಲೋಕದಲ್ಲಿಯೂ ಸಹ ತಮ್ಮ ಗುರುಗಳ ಅಥವಾ ದೈವದ ರಕ್ಷಣೆ ಸದಾ ಇರುತ್ತದೆ.
\end{enumerate}

ಕುರಂಟಕ ಪದ್ಧತಿಯು ಇಲ್ಲಿಗೆ ಸಂಪೂರ್ಣವಾಯಿತು.
%ಕೊರಂಟಕ ಪದ್ಧತಿಯ ತೃತೀಯ ಅಧ್ಯಾಯವು ಇಲ್ಲಿಗೆ ಮುಕ್ತಾಯವಾಯಿತು.
\end{kan}
\end{multicols}
